% Slajd 4 – mapa demonstracji na żywo (pkt 6 wyzwania: określić potrzeby grupy, sprawdzić miejsce,
% pokazać konkretne bariery i udogodnienia, źródło i datę informacji, oznaczenie danych niepełnych
% i przykładowych). Scenariusz kończy się sukcesem: Anna znajduje miejsce, które spełnia wszystkie jej potrzeby.
% Slajd wisi na ekranie tuż przed przejściem do aplikacji; szczegóły w notatkach.
\begin{frame}{Demo: Anna na wózku znajduje restaurację bez barier}
  \begin{columns}[T,onlytextwidth]
    \begin{column}{0.63\textwidth}
      \scriptsize
      \renewcommand{\arraystretch}{1.1}
      \begin{tabularx}{\textwidth}{@{}>{\bfseries\color{accNavy}}c >{\raggedright\arraybackslash}p{30mm} >{\raggedright\arraybackslash}X@{}}
        \toprule
        & \textbf{Krok} & \textbf{Na co zwrócić uwagę} \\
        \midrule
        1 & Preferencje: bez schodów, toaleta dostępna, parking OzN &
            mapa i~listy pokazują tylko bariery ważne dla Anny; miejsca posortowane według dopasowania \\
        2 & Pytanie: „restauracja na Kazimierzu bez schodów, z~toaletą i~parkingiem” &
            ranking z~uzasadnieniem: Kuchnia u~Doroty na górze spełnia wszystko; przy innych widać, czego \emph{nie wiadomo} \\
        3 & Karta miejsca: Kuchnia u~Doroty &
            zielone „Spełnia wszystkie Twoje potrzeby”; każdy fakt ze źródłem, datą i~statusem; braki nazwane wprost \\
        4 & Po drodze: zgłoszenie „Brak windy” na przystanku, waga „blokuje” &
            kwadrat = blokuje; głosy „nadal jest” / „naprawione”; awaria wygasa po 14~dniach \\
        5 & Panel właściciela: „Potwierdź informacje” &
            status „Potwierdzone przez właściciela”; Anna dostaje powiadomienie \\
        \bottomrule
      \end{tabularx}

      \vspace{1.5mm}
      \textbf{Przypadki brzegowe w~demo:} dane sprzeczne $\rightarrow$ „Wymaga ponownej weryfikacji”;
      dane miasta niedostępne $\rightarrow$ snapshot z~datą „stan na”; brak informacji nigdy nie znaczy „dostępne”.

      \vspace{1mm}
      \muted{Dane przykładowe są oznaczone: konta demo (\texttt{anna@}, \texttt{kasia@}, \texttt{admin@accessly.test}),
        zgłoszenia testowe, etykieta „Tryb demo”.}
    \end{column}
    \begin{column}{0.35\textwidth}
      \screenshot{asystent-ok-wynik}\par
      \vspace{0.3mm}
      \muted{Krok 2: miejsce, które spełnia wszystko -- zielone znaczniki i~„Dopasowane do Ciebie”.}
      \vspace{1.5mm}
      \screenshot{karta-ok-atrybut}\par
      \vspace{0.3mm}
      \muted{Krok 3: jeden fakt -- wartość, źródło, data, status i~przyciski Aktualne / Nieaktualne / Popraw.}
    \end{column}
  \end{columns}

  \note[item]{Czas: ok. 3 minuty. Przed wejściem: aplikacja na telefonie i~laptopie, mapa na Kazimierzu, w~Profilu zaznaczone preferencje bez schodów, toaleta dostępna, parking OzN (bez chipa „Wózek”: wymaga on też szerokości drzwi, której nie ma w~danych OSM, więc dawałby „Spełnia 2~z~3”), panel właścicielki otwarty w~drugiej karcie.}
  \note[item]{Krok 1 (0:00–0:30): pokazać Profil $\rightarrow$ „Moje preferencje” i~listę „Najbliżej środka mapy”; miejsca z~odznaką „Dopasowane do Ciebie” są wyżej.}
  \note[item]{Krok 2 (0:30–1:00): zakładka Miejsca, wpisać „Restauracja na Kazimierzu bez schodów, z~dostępną toaletą i~parkingiem” – w~trybie reguł wynik jest deterministyczny: „Znalazłem 535 pasujących miejsc; 1~z~pokazanych spełnia wszystkie Twoje wymagania”, na górze Kuchnia u~Doroty.}
  \note[item]{Krok 3 (1:00–1:45): otworzyć kartę; zielona ramka; pokazać palcem trzy rzeczy przy atrybucie: źródło, datę, status. Przewinąć do atrybutu bez danych (np. szerokie drzwi) – nasz najważniejszy komunikat: nie zgadujemy. Przyciski Aktualne / Nieaktualne / Popraw, „Uzupełnij” i~pytanie do właściciela.}
  \note[item]{Krok 4 (1:45–2:20): „Zgłoś problem”, pinezka na przystanku Stradom, typ „Brak windy”, waga „blokuje”; pokazać kształt znacznika (kwadrat = blokuje) i~głosy. Zgłoszenie można zrobić bez konta.}
  \note[item]{Krok 5 (2:20–3:00): w~drugiej karcie \texttt{/owner.html} jako \texttt{kasia@accessly.test} (Camelot Cafe) kliknąć „Potwierdź informacje”; wrócić do karty tego lokalu – status się zmienił.}
  \note[item]{Przypadki brzegowe (gdy jury pyta lub zostaje czas): chip „Wózek” i~„Spełnia 2~z~3, brak danych: szerokie drzwi”; historia atrybutu ze sprzecznymi głosami; strona „Źródła danych” z~datami snapshotów.}
  \note[item]{Jeśli padnie sieć: wszystko poza kafelkami mapy działa z~danych lokalnych. Dane przykładowe: trzy konta demo i~zgłoszenia testowe; widok administratora bez logowania ma etykietę „Tryb demo: dane przykładowe”.}
\end{frame}
